\section{Discussion}
\label{sec:discussion}

The introduction of CoShard as a representation-space partitioning methodology offers a pragmatic approach for managing privacy-utility trade-offs in distributed RAG architectures. In this section, we discuss the broader implications of our findings across system design, practical risk reduction, and deployment constraints.

\subsection{Practical Risk Reduction in RAG Systems}
The rapid integration of LLMs into enterprise workflows often relies on centralized vector databases, which pool an organization's intellectual property into a single attack surface. Under strict data governance frameworks, organizations are encouraged to follow principles of \textit{Data Minimization} and \textit{Compartmentalization}.

Semantic sharding, guided by the CoShard optimization objective, provides a quantifiable mechanism for progressing toward these principles. By explicitly enforcing a low SHADE penalty, CoShard acts as a structural deterrent against single-shard concentration. If a data breach occurs—e.g., an insider threat exfiltrating a specific shard—the adversary receives a highly fragmented subset of the corpus. While this does not provide formal cryptographic privacy guarantees (such as those offered by Fully Homomorphic Encryption or Differential Privacy), it significantly reduces the empirical probability that an adversary can reconstruct a contiguous narrative regarding a specific sensitive entity, thereby limiting the practical blast radius of a localized breach.

\subsection{Technical and Systems Design}
From a distributed systems perspective, the empirical privacy-utility trade-off dictates a strict architectural decision. Modularity-maximizing clustering (e.g., Leiden) is conceptually straightforward but disregards data concentration. The CoShard pipeline requires the construction of an explicit K-NN similarity graph, which naively scales at $O(n^2)$. 

For massive enterprise corpora containing millions of dense embeddings, constructing the exact K-NN graph becomes computationally intensive. However, CoShard can be readily adapted to operate over Approximate Nearest Neighbor (ANN) graphs, such as Hierarchical Navigable Small World (HNSW) graphs, which reduce the construction complexity to $O(n \log n)$. Furthermore, the 2-phase modularity maximization algorithm employed by CoShard is parallelizable. The graph contraction step rapidly reduces the search space, allowing CoShard to scale logarithmically after the initial graph construction phase. 

Future systems engineering research must explore the dynamic updating of the CoShard graph. As new documents are continuously ingested into the RAG pipeline, the system must perform localized, online re-sharding to maintain the Pareto optimal boundary without requiring a full re-computation of the global SHADE metric.

\subsection{Deployment Constraints and Broader Impact}
The current gold standard for privacy-preserving ML relies heavily on Federated Learning and Confidential Computing (e.g., Secure Enclaves). While these architectures provide rigorous cryptographic security, they require significant capital expenditure, specialized hardware, and complex key management infrastructures.

CoShard offers a complementary, software-layer approach by shifting the initial privacy burden to intelligent data topology. By relying on algorithmic graph partitioning to enforce privacy via information starvation, smaller organizations can deploy risk-aware RAG agents using commodity hardware. 

It is crucial to note that CoShard is not a replacement for traditional access controls or endpoint security. Rather, it serves as a defense-in-depth mechanism. By structurally starving potential adversaries of cohesive semantic context, CoShard elevates the difficulty of embedding inversion attacks, providing a necessary, scalable safeguard for privacy-aware generative AI deployments.
